Finally, we are explicit about what the model does not settle. It takes no stance on
\emph{why} an informed seller withholds, since deliberate concealment and an unrecorded but
still-valid certificate are observationally similar in the listing data, and we do not
attempt to separate them. It assumes no falsification, so a displayed score is taken at face
value. And it maintains Corollary~\ref{cor:hq_disclose}, $\sigma_1^*=0$, rather than testing
it; Section~\ref{ss:attention} relaxes the buyer side and shows where this restriction survives. Informed non-disclosers are assigned low quality by revealed preference, which is an
assumption and not a result. Appendix~\ref{app:robustness} reports estimates that instead read
the withheld quality from the registry, so the restriction can be assessed. Nor does the
design pin down $\rho_m$ from the interior in most markets, where the data are consistent with full
buyer sensitivity but cannot establish how close to one it lies, and we maintain
$\hat\rho_m = 1$ for those markets and treat $\hat c_{0,m}$ as an upper bound accordingly. We
do not purport to estimate the one true model of housing disclosure; we propose one way of
reading these data through a workhorse framework.
